% Historical source (removed by author request): references/_OLD_Übungsaufgaben/Übungsblatt 3 - Vektorechnung/Blatt3.tex:69
\begin{ExerciseBlock}
\ExerciseHeading{04}{BAC-CAB-Regel}{vector_algebra_bac_cab_identity}
Für das doppelte Kreuzprodukt gilt die sogenannte \textit{bac-cab-Regel}.

\begin{equation*}
	\vec{a} \times \left( \vec{b} \times \vec{c} \right)
	=
	\vec{b} \left( \vec{a} \cdot \vec{c} \right)
	-
	\vec{c} \left( \vec{a} \cdot \vec{b} \right)
\end{equation*}

Rechnen Sie die Richtigkeit der Regel allgemein nach. Hinweis: Verwenden Sie eine allgemeine Koordinatendarstellung von 3 Vektoren in 3D: \par
\begin{center} 
$ \vec{a}=\left(\begin{array}{c} a_1 \\ a_2 \\ a_3 \end{array}\right)$, 
$\vec{b}=\left(\begin{array}{c} b_1 \\ b_2 \\ b_3 \end{array}\right)$ und ${\vec{c}}=\left(\begin{array}{c} c_1 \\ c_2 \\ c_3 \end{array}\right)$. 
\end{center}

\end{ExerciseBlock}
\UseExerciseSolution{vector_algebra_bac_cab_identity}
